\section[Results]{The results}
\sectionframe{03}{The results}{Does it work?}{\faChartBar}

% ------------------------------------------------------------------ the error
% One plot per message: the method in the primary colour, the others in grey, the number that matters
% marked in the accent colour.
\begin{frame}[source={Author et al.\ (2026), Journal 1, 1, Table 2}]{Error on the test set}
  \begin{tikzpicture}[x=1mm,y=1mm]
    \useasboundingbox (0,0) rectangle (144,-54);
    % axes: 80 mm per unit of error, from 0 at y = -46
    \draw[greymid,line width=0.5pt] (12,-2) -- (12,-46) -- (96,-46);
    \foreach \v in {0,0.1,0.2,0.3,0.4,0.5} {
      \draw[greymid,line width=0.5pt] (12,{-46+80*\v}) -- ++(-1,0);
      \node[anchor=east,small label] at (10.2,{-46+80*\v}) {\v};}
    \node[anchor=south,small label,rotate=90] at (2.4,-26) {error};
    % bars, 12 mm wide: left edge, value, colour, label
    \foreach \x/\v/\c/\l in {18/0.42/greymid/{$\lambda=0$, baseline},46/0.27/primary/{$\lambda=0.1$},
        74/0.33/greymid/{$\lambda=1$}} {
      \fill[\c] (\x,-46) rectangle ++(12,{80*\v});
      \node[anchor=south,label text] at ({\x+6},{-46+80*\v+0.8}) {\v};
      \node[anchor=north,small label,text=ink] at ({\x+6},-47.2) {\l};}
    % the improvement, beside the method's bar
    \draw[greymid,line width=0.4pt,dash pattern=on 1.2pt off 1pt] (30,{-46+80*0.42}) -- (59.6,{-46+80*0.42});
    \draw[accent,line width=0.6pt] (59.6,{-46+80*0.42}) -- ++(1.6,0) -- (61.2,{-46+80*0.27}) -- ++(-1.6,0);
    \node[anchor=west,label text,text=accent,font=\footnotesize\bfseries] at (62.4,{-46+80*0.345}) {−36\,\%};
    \node[anchor=north west,small label,text width=40mm] at (102,-6)
      {mean absolute error over the test images; lower is better};
  \end{tikzpicture}
  \note{\textbf{6:10--7:10}\par
    The mean absolute error on the test images, lower is better.\par
    With a weight of 0.1 the error falls from 0.42 to 0.27, by about a third.\par
    A weight of 1 is worse than 0.1: the halos of the previous slide cost pixel error.\par
    These are example values; put your own results and their spread here.}
\end{frame}

% ------------------------------------------------------- the case, revisited
% The failure case of section 1 returns, now with the method's output beside it.
\begin{frame}[source={Author et al.\ (2026), Journal 1, 1, Fig.\ 2}]{Baseline against method}
  \begin{tikzpicture}[x=1mm,y=1mm]
    \useasboundingbox (0,0) rectangle (144,-54);
    \foreach \k/\p/\l in {0/0/input,1/1000/baseline,2/2000/{method, $\lambda=0.1$}} {
      \figpart{16+\k*39}{0}{assets/example-figure}{4000}{\p}{0}{\p+1000}{1000}{34}
      \node[small label,anchor=north] at ({33+\k*39},-35.2) {\l};}
    \goodstrip{-40.6}{the thin bar and the shared edge stay sharp}
    \warnstrip{-48.4}{one dataset and one training run so far}
  \end{tikzpicture}
  \note{\textbf{7:10--8:10}\par
    The case from the beginning, with the method's output on the right.\par
    The thin bar and the edge between the two squares stay sharp.\par
    The evidence is one dataset and one training run; the spread over several runs is the next
      check.}
\end{frame}
